% Source printed pp. 204-209 / PDF pp. 212-217.
\section{Deligne--Mumford 叠的有限覆盖}\label{src:5.6}
\begin{epigraph}
想一想就显而易见了.
\epigraphauthor{Angelo Vistoli}
\end{epigraph}
我们证明 Gabber 引理 (\bookref{5.6.1}): 分离 Deligne--Mumford 叠有来自概形的有限覆盖. 作为推论, 得到代数空间 (\bookref{5.6.6}) 与 Deligne--Mumford 叠 (\bookref{5.6.5}) 的 Chow 引理. 在~\S\bookref{5.6.3-subsection} 中推导其上同调应用.

\subsection{Gabber 引理}\label{src:5.6.1-subsection}
一个归功于 Gabber 的著名结果断言, Deligne--Mumford 叠有来自概形的有限覆盖. 论证最初发表于~\cite{Del85,Vis89}.
\begin{theorem}[Gabber 引理]\label{src:5.6.1}
设 $\mathcal X$ 是 Noether 概形 $S$ 上的分离有限型 Deligne--Mumford 叠. 则存在来自概形 $Z$ 的有限, 一般平展且满射的态射 $Z\to\mathcal X$.
\end{theorem}
\begin{proof}
将 $\mathcal X$ 换成具有约化叠结构的不可约分支的不交并后, 可假设 $\mathcal X$ 不可约且约化. 每个平展呈示 $U\to\mathcal X$ 都可表出, 拟有限且分离, 因而由 Zariski 主定理 (\bookref{5.5.9}), 它可分解为开浸入 $U\hookrightarrow\widetilde{\mathcal X}$ 与有限态射 $\widetilde{\mathcal X}\to\mathcal X$ 的复合. 将 $\mathcal X$ 换成 $\widetilde{\mathcal X}$ 后, 可假设 $\mathcal X$ 有稠密开子概形. 若 $p:U\to\mathcal X$ 是平展呈示, 则存在稠密开子概形 $V\subseteq\mathcal X$, 使得 $p^{-1}(V)\to V$ 是次数为 $d$ 的有限平展态射. 可选取有限平展覆盖 $V'\to V$, 使得 $p^{-1}(V)\times_V V'\to V'$ 是平凡平展覆盖; 事实上, 如命题~\bookref{A.3.11}, 可取 $V'$ 为
\[
(V'/V)^d=\underbrace{V'\times_V\cdots\times_V V'}_{d}
\]
中两两对角线 (开且闭) 的补集. 再次应用 Zariski 主定理 (\bookref{5.5.9}), 这次应用于复合 $V'\twoheadrightarrow V\hookrightarrow\mathcal X$, 得到有限满射态射 $\widetilde{\mathcal X}\to\mathcal X$, 其限制为 $V'\to V$. 因而将 $\mathcal X$ 换成 $\widetilde{\mathcal X}$ 后, 还可假设存在平展呈示 $U\to\mathcal X$, 它在稠密开子概形 $j:V\hookrightarrow\mathcal X$ 上是平凡平展覆盖, 即有 Cartesian 图
\[
\begin{tikzcd}
\coprod_{i=1}^d V\arrow[r,hook,"{j'}"]\arrow[d]&U\arrow[d,"p"]\\
V\arrow[r,hook,"j"']&\mathcal X.
\end{tikzcd}
\]
将构造来自概形的有限满射态射 $Z\to\mathcal X$, 它在 $V$ 上是同构. 设 $\mathcal A\subseteq j_*\mathcal O_Z$ 是 $\mathcal O_{\mathcal X}\to j_*\mathcal O_Z$ 的整闭包. 则 $p^*\mathcal A$ 是 $\mathcal O_U$ 在 $j'_*\mathcal O_{\coprod V}=p^*j_*\mathcal O_V$ 中的整闭包 (命题~\bookref{A.7.4}). 定义第 $i$ 个 $V$ 副本的幂等元 $e_i\in\Gamma(U,j'_*\mathcal O_{\coprod V})=\Gamma(\coprod V,\mathcal O_{\coprod V})$ 在 $\mathcal O_U$ 上整, 因而定义整体截面 $e_i\in\Gamma(U,p^*\mathcal A)$. 现在将 $\mathcal A=\operatorname{colim}_\lambda\mathcal A_\lambda$ 写成有限型 $\mathcal O_{\mathcal X}$-子代数的滤过余极限 (习题~\bookref{5.1.33}). 由于 $\mathcal A$ 在 $\mathcal O_{\mathcal X}$ 上整, 每个 $\mathcal A_\lambda$ 都是有限 $\mathcal O_{\mathcal X}$-代数. 对 $\lambda\gg0$, 有 $e_i\in\Gamma(U,p^*\mathcal A_\lambda)$. Deligne--Mumford 叠 $Z:=\Spec_{\mathcal X}\mathcal A_\lambda$ 在 $\mathcal X$ 上有限, 断言 $Z$ 是概形. 为此, 考虑 Cartesian 图
\[
\begin{tikzcd}
Z'\arrow[r]\arrow[d]&U\arrow[d,"p"]\\
Z\arrow[r]&\mathcal X\text{\textit{cartesian}},
\end{tikzcd}
\]
其中 $Z'$ 是概形 (因为它在 $U$ 上有限). 每个幂等元 $e_i$ 定义 $\mathcal O_{Z'}$ 的整体截面, 因而给出分解 $Z'=\coprod_{i=1}^d Z'_i$. 每个态射 $Z'_i\to Z$ 都平展, 分离且双有理, 因而是开浸入. 由于 $Z'\to Z$ 满射, $\{Z'_i\}$ 定义 $Z$ 的开覆盖, $Z$ 是概形. 另见~\cite[定理 $16.6$]{LMB00}.
\end{proof}
\begin{exercise}\label{src:5.6.2}
仿照~\cite[命题 $2.6$]{Vis89} 给出另一证明:
\begin{enumerate}[(i)]
\item 用极限方法归结到 $S$ 在 $\mathbb Z$ 上有限型的情形 (这保证正规化有限).
\item 归结到 $\mathcal X$ 不可约, 约化且正规的情形.
\item 设 $\mathcal X\to X$ 是粗模空间, $U\to\mathcal X$ 是来自概形 $U=\coprod_i U_i$ 的平展呈示, 其中每个 $U_i$ 都整且仿射. 设 $L$ 是 $K(X)$ 的正规有限域扩张, 包含每个函数域 $K(U_i)$, $Y\to X$ 与 $Y_i\to U_i$ 分别为 $X$ 与 $U_i$ 在 $L$ 中的正规化. 证明 $Y$ 是概形, 且 $Y\to X$ 在 Zariski 局部经 $\mathcal X$ 分解: 说明每个 $Y_i\to Y$ 都是开浸入, 并且 $Y$ 是所有 $i$ 及 $\sigma\in\operatorname{Aut}(L/K(X))$ 下平移 $\sigma\cdot Y_i$ 的并.
\item 将 $X$ 换成 $Y$, $\mathcal X$ 换成 $\mathcal X\times_X Y$ 后, 归结到粗模空间 $X$ 是概形且有开覆盖 $X=\bigcup X_\alpha$ 的情形, 其中每个 $X_\alpha\hookrightarrow X$ 都经 $\mathcal X$ 分解. 证明将 $X$ 换成有限覆盖后, 截面 $X_\alpha\to\mathcal X$ 可粘合为整体截面 $X\to\mathcal X$.
\end{enumerate}
\end{exercise}
\begin{remark}\label{src:5.6.3}
更一般地, $S$ 上的每个分离有限型代数叠都有来自拟射影概形的固有覆盖~\cite{Ols05}.
\end{remark}
\begin{exercise}[很好的练习]\label{src:5.6.4}
设 $X$ 是 Noether 概形 $S$ 上的正规有限型代数空间. 证明存在正规概形 $U$ 及有限抽象群 $G$ 的作用, 使得 $X$ 是 $U$ 关于 $G$ 的商, 即 $[U/G]\to X$ 是粗模空间.

\textit{提示:} 归结到 $X$ 整的情形后, 选取来自概形的有限, 一般平展且满射的态射 $U\to X$. 设 $L$ 是 $K(U)/K(X)$ 的 Galois 闭包, 取 $U$ 为 $X$ 在 $L$ 中的整闭包, 取 $G=\operatorname{Gal}(L/K(X))$. 另见~\cite[推论 $16.6.2$]{LMB00}.
\end{exercise}

\subsection{Chow 引理}\label{src:5.6.2-subsection}
\begin{corollary}[Deligne--Mumford 叠的 Chow 引理]\label{src:5.6.5}
设 $\mathcal X$ 是 Noether 概形 $S$ 上的分离有限型 Deligne--Mumford 叠. 则存在来自 $S$ 上拟射影概形 $Z$ 的射影, 一般平展且满射的态射 $Z\to\mathcal X$.
\end{corollary}
\begin{proof}
Gabber 引理 (\bookref{5.6.1}) 将此陈述归结到概形的情形, 参见~\cite[习题 II.4.10]{Har77}.
\end{proof}
代数空间还有一个双有理版本.
\begin{exercise}[代数空间的 Chow 引理]\label{src:5.6.6}
设 $\mathcal X$ 是 Noether 概形 $S$ 上的分离有限型代数空间. 则存在来自 $S$ 上拟射影概形 $Z$ 的射影, 双有理且满射的态射 $Z\to\mathcal X$. 另见~\cite[IV.3.1]{Knu71} 和~\cite[Tag 088U]{SP}.
\end{exercise}
\begin{exercise}[赋值判据可在稠密开集上检验]\label{src:5.6.7}
\begin{enumerate}[(a)]
\item 设 Noether 概形 $S$ 上的 Deligne--Mumford 叠 $\mathcal X$ 分离且有限型, $\mathcal U\subseteq\mathcal X$ 是稠密开子叠. 证明 $\mathcal X\to S$ 固有, 当且仅当对 $S$ 上定义的每个离散赋值环 $R$, 设其分式域为 $K$, 将 $R$ 换成离散赋值环的扩张后, 每个态射 $\Spec K\to\mathcal U$ 都可延拓为 $S$ 上的态射 $\Spec R\to\mathcal X$.

\textit{提示:} 用推论~\bookref{5.6.5} 归结到 $\mathcal X\to S$ 拟射影的情形.
\item 设 $\mathcal X$ 是对角线分离且拟紧的 Deligne--Mumford 叠. 证明到 Noether 概形的有限型态射 $\mathcal X\to S$ 分离, 当且仅当对 $S$ 上的每个离散赋值环 $R$, 设其分式域为 $K$, 以及每对 $S$ 上的态射 $h,g:\Spec R\to\mathcal X$, 其一般纤维的任意同构 $h_K\xrightarrow{\sim}g_K$ 都可延拓为唯一同构 $h\xrightarrow{\sim}g$.

\textit{提示:} 归结到 (a).
\end{enumerate}
另见~\cite[II.7.3.10]{EGA}, \cite{MO09} 和~\cite[Tag 089G]{SP}.
\end{exercise}

\subsection{对上同调的应用}\label{src:5.6.3-subsection}
固有概形上凝聚层上同调的大多数基本定理---上同调有限性 (\bookref{A.5.3}), 上同调与基变换 (\bookref{A.6.2}), 半连续性 (\bookref{A.6.4}), 形式函数 (\bookref{A.5.4}), Grothendieck 存在定理 (\bookref{C.5.3}) 及 Stein 分解~\cite[推论 III.11.5]{Har77}---对固有 Deligne--Mumford 叠也成立. 它们蕴含 Zariski 连通性定理 (\bookref{5.6.15}), 这是证明正特征下 $\mathcal M_g$ 不可约性 (定理~\bookref{6.7.33}) 的一个关键工具.
\begin{theorem}[上同调有限性]\label{src:5.6.8}
设 $f:\mathcal X\to\mathcal Y$ 是 Noether Deligne--Mumford 叠的固有态射. 对 $\mathcal X$ 上的任意凝聚层 $F$ 和任意 $i\geq0$, $R^if_*F$ 凝聚.
\end{theorem}
\begin{proof}
由平坦基变换 (\bookref{5.1.43}), 可假设 $\mathcal Y=\Spec A$ 是仿射概形, 须说明对 $i\geq0$, $\mathrm H^i(\mathcal X,F)$ 是有限 $A$-模. 将使用\newterm{逐层分解}{dévissage} (\bookref{A.5.1}) 的以下推广, 它的证明相同: 若 $\mathcal P$ 是 $\mathcal X$ 上凝聚层的性质, 满足
\begin{enumerate}[(a)]
\item 在 $\operatorname{Coh}(\mathcal X)$ 的短正合列中, 若三个对象中的两个满足 $\mathcal P$, 则第三个也满足, 且
\item 对具有理想层 $I_{\mathcal Z}\subseteq\mathcal O_{\mathcal X}$ 的所有整闭子叠 $\mathcal Z\subseteq\mathcal X$, 以及 $\mathcal X$ 上满足 $I_{\mathcal Z}F=0$ 的凝聚层 $F$, 存在 $\mathcal X$ 上满足 $\mathcal P$ 的凝聚层 $G$, 使得 $I_{\mathcal Z}G=0$, 并有态射 $F\to G$, 它在 $\mathcal Z$ 的某个开子叠上是同构,
\end{enumerate}
则 $\mathcal X$ 上的每个凝聚层都满足 $\mathcal P$. 取 $\mathcal P$ 为对 $i\geq0$, $\mathrm H^i(\mathcal X,F)$ 是有限 $A$-模这一性质, 可知 (a) 成立. 对 (b), 用 Gabber 引理 (\bookref{5.6.1}), 取来自概形的有限满射态射 $g_0:V_0\to\mathcal Z$. 置 $V_1=V_0\times_{\mathcal Z}V_0$, 其投影为 $g_1:V_1\to\mathcal Z$, 有复形 $0\to F\to g_{0,*}g_0^*F\to g_{1,*}g_1^*F$, 当 $V_0\to\mathcal Z$ 平坦时它正合 (命题~\bookref{3.1.1}). 定义 $G:=\ker(g_{0,*}g_0^*F\to g_{1,*}g_1^*F)$, 则有映射 $F\to G$. 由一般平坦性 (\bookref{4.3.32}), $V_0\to\mathcal Z$ 在非空开子叠 $\mathcal U\subseteq\mathcal Z$ 上平坦, 从而 $F\to G$ 在 $\mathcal U$ 上是同构. 由于 $V_0$ 与 $V_1$ 都在 $\Spec A$ 上固有, $g_{0,*}g_0^*F$ 与 $g_{1,*}g_1^*F$ 有有限上同调, 因而 $G$ 也如此. 另见~\cite[定理 $1$]{Fal03}, \cite[定理 $1.2$]{Ols05}, \cite[定理 $15.6$]{LMB00} 和~\cite[定理 $11.6.1$]{Ols16}.
\end{proof}
\begin{remark}\label{src:5.6.9}
与概形不同, $\mathrm H^i(\mathcal X,F)$ 可对任意大的 $i$ 非零; 见习题~\bookref{5.1.47}.
\end{remark}
\begin{exercise}\label{src:5.6.10}
若 $X$ 是 Noether 代数空间, 证明存在整数 $N$, 使得对 $X$ 上的每个凝聚层 $F$, 当 $i>N$ 时都有 $\mathrm H^i(X,F)=0$.
\end{exercise}
\begin{proposition}[半连续性]\label{src:5.6.11}
设 $\mathcal X\to\mathcal Y$ 是 Noether Deligne--Mumford 叠的固有态射, $F$ 是 $\mathcal X$ 上在 $\mathcal Y$ 上平坦的凝聚层. 对每个 $i\geq0$, 函数
\begin{equation}\label{src:5.6.12}
|\mathcal Y|\longrightarrow\mathbb Z,\qquad y\longmapsto\dim_K\mathrm H^i(\mathcal X_K,F_K),
\end{equation}
上半连续, 其中 $\Spec K\to\mathcal Y$ 是代表 $y$ 的任意域值点.
\end{proposition}
\begin{proof}
留给读者检验函数~\bookeqref{5.6.12} 可构造. 因而只需检验它关于 $|\mathcal Y|$ 中的特殊化 $y\rightsquigarrow y_0$ 上半连续. 由命题~\bookref{4.8.21}, 此特殊化可由来自离散赋值环的态射 $\Spec R\to\mathcal Y$ 实现, 因而可假设 $\mathcal Y=\Spec R$. 设 $K$ 与 $\kappa$ 分别为 $R$ 的分式域与剩余域, 须说明 $\dim_K\mathrm H^i(\mathcal X_K,F_K)\leq\dim_\kappa\mathrm H^i(\mathcal X_\kappa,F_\kappa)$. 若 $\pi\in R$ 是一致化元, 则由于 $F$ 在 $R$ 上平坦, $\pi:F\to F$ 单射. 对短正合列 $0\to F\to F\to F_\kappa\to0$ 应用 $\mathrm H^i(\mathcal X,-)$, 得到正合列
\[
\cdots\longrightarrow\mathrm H^i(\mathcal X,F)\xrightarrow{\pi}\mathrm H^i(\mathcal X,F)\longrightarrow\mathrm H^i(\mathcal X_\kappa,F_\kappa)\longrightarrow\cdots.
\]
这给出单射 $\mathrm H^i(\mathcal X,F)/\pi\mathrm H^i(\mathcal X,F)\hookrightarrow\mathrm H^i(\mathcal X_\kappa,F_\kappa)$. 另一方面, 由上同调有限性 (\bookref{5.6.8}), $\mathrm H^i(\mathcal X,F)$ 是有限 $R$-模, 由平坦基变换 (\bookref{5.1.43}) 有 $\mathrm H^i(\mathcal X,F)\otimes_R K=\mathrm H^i(\mathcal X,F_K)$. 这给出
\[
\dim_K\mathrm H^i(\mathcal X_K,F_K)\leq\dim_\kappa\mathrm H^i(\mathcal X,F)/\pi\mathrm H^i(\mathcal X,F)\leq\dim_\kappa\mathrm H^i(\mathcal X_\kappa,F_\kappa).
\]
\end{proof}
\begin{exercise}[困难]\label{src:5.6.13}
设 $\mathcal X$ 是完备 Noether 局部环 $(A,\mathfrak m)$ 上的固有 Noether Deligne--Mumford 叠. 置 $\mathcal X_n=\mathcal X\times_A A/\mathfrak m^{n+1}$.
\begin{enumerate}[(a)]
\item (形式函数) 若 $F$ 是 $\mathcal X$ 上的凝聚层, 则对每个 $i\geq0$, 存在自然同构
\[
\mathrm H^i(\mathcal X,F)\xrightarrow{\sim}\varprojlim_n\mathrm H^i(\mathcal X_n,F|_{\mathcal X_n}).
\]
\item (Grothendieck 存在定理) 存在范畴等价
\[
\operatorname{Coh}(\mathcal X)\longrightarrow\varprojlim_n\operatorname{Coh}(\mathcal X_n),\qquad F\longmapsto\{F|_{\mathcal X_n}\}.
\]
\end{enumerate}
另见~\cite[定理 $1.4$]{Ols05} 和~\cite[推论 $3.2$, 定理 $4.1$]{Con05a}.
\end{exercise}
\begin{theorem}[Stein 分解/Zariski 连通性定理]\label{src:5.6.14}
Noether Deligne--Mumford 叠的固有态射 $f:\mathcal X\to\mathcal Y$ 可分解为
\[
f:\mathcal X\xrightarrow{f'}\mathcal Y'=\Spec_{\mathcal Y}f_*\mathcal O_{\mathcal X}\xrightarrow{g}\mathcal Y,
\]
其中 $f'$ 有几何连通的纤维, $g$ 有限. 特别地, 若 $f_*\mathcal O_{\mathcal X}=\mathcal O_{\mathcal Y}$, 则 $f$ 有几何连通的纤维.
\end{theorem}
\begin{proof}
这与概形情形的证明完全相同---\cite[推论 III.11.5]{Har77}, \cite[III$_1$.4.3.1]{EGA} 和~\cite[Tag 03H2]{SP}---使用形式函数. 由上同调有限性 (\bookref{5.6.8}), $f_*\mathcal O_{\mathcal X}$ 凝聚, $g$ 有限. 因而归结到证明若 $f_*\mathcal O_{\mathcal X}=\mathcal O_{\mathcal Y}$, 则 $f$ 有几何连通的纤维. 由于问题是平展局部的, 还可假设 $\mathcal Y=S$ 仿射. 此外, 由于纤维 $\mathcal X_s$ 的几何连通性可在 $\kappa(s)$ 的可分有限域扩张上检验, 只需验证对每个平展态射 $S'\to S$, $f':\mathcal X_{S'}\to S'$ 的纤维连通. 由平坦基变换 (\bookref{5.1.43}), $f'_*\mathcal O_{\mathcal X_{S'}}=\mathcal O_{S'}$, 因而只需验证 $f$ 有连通纤维. 若纤维 $\mathcal X_s$ 不连通, 置 $S_n=\Spec\mathcal O_{S,s}/\mathfrak m_s^{n+1}$, 则对非零凝聚 $\mathcal O_{S_n}$-模 $A_n$ 与 $B_n$, 存在相容的同构 $\mathrm H^0(\mathcal X\times_S S_n,\mathcal O_{\mathcal X\times_S S_n})\cong A_n\times B_n$. 由形式函数 (\bookref{5.6.13}(b)), $\widehat{\mathcal O}_{S,s}\cong\varprojlim A_n\times\varprojlim B_n$, 这与 $\widehat{\mathcal O}_{S,s}$ 是局部环矛盾.
\end{proof}
以下结果将应用于态射 $\mathcal M_g\to\Spec\mathbb Z$, 以说明 $\mathcal M_g\times_{\mathbb Z}\mathbb F_p$ 连通 (定理~\bookref{6.7.33}). 此应用只需要下面的下半连续性, 它依赖 Stein 分解 (\bookref{5.6.14}), 而不依赖半连续性 (\bookref{5.6.11}).
\begin{corollary}[Zariski 连通性定理 II]\label{src:5.6.15}
设 $f:\mathcal X\to\mathcal Y$ 是 Noether Deligne--Mumford 叠的固有平坦态射. 函数
\begin{equation}\label{src:5.6.16}
|\mathcal Y|\longrightarrow\mathbb Z,\qquad y\longmapsto\#\{\mathcal X\times_{\mathcal Y}\Spec K\text{ 的连通分支}\},
\end{equation}
下半连续, 其中 $\Spec K\to\mathcal Y$ 是代表 $y$ 的任意几何点. 若此外 $f$ 有几何约化的纤维, 则~\bookeqref{5.6.16} 局部常值.
\end{corollary}
\begin{proof}
由平展下降, 可假设 $\mathcal Y=S$ 是概形. 设 $\mathcal X\xrightarrow{f'}S'\to S$ 是 Stein 分解 (\bookref{5.6.14}). 由于 $S'\to S$ 有限, 点 $s\in S$ 上方只有有限多个点 $s'_1,\ldots,s'_m$. 由拟有限态射的平展局部化 (\bookref{A.7.1}), 将 $S$ 换成 $s$ 的平展邻域后, 可使 $S'=S'_1\amalg\cdots\amalg S'_m$, 其中 $s'_i\in S'_i$, 且 $\kappa(s'_i)/\kappa(s)$ 纯不可分. 由于 $\mathcal X\to S'$ 有几何连通的纤维, 每个纤维 $\mathcal X_{s'_i}$ 都几何连通, 因而 $\mathcal X\to S$ 在 $s$ 上的几何纤维恰有 $m$ 个连通分支. 由于 $\mathcal X\to S$ 平坦, $f'^{-1}(S'_i)$ 在 $S$ 中的像 $U_i\subseteq S$ 是开的. 在 $s$ 的开邻域 $\bigcap_i U_i$ 中, 每个几何纤维至少有 $m$ 个连通分支.

现在假设每个纤维都几何约化, 则 $s\in S$ 上的几何分支数恰为 $\dim_{\kappa(s)}\mathrm H^0(\mathcal X_s,\mathcal O_{\mathcal X_s})$. 由半连续性 (\bookref{5.6.11}), \bookeqref{5.6.16} 也上半连续. 概形的情形另见~\cite[IV$_3$.15.5.3--7]{EGA} 和~\cite[Tags 0BUI 和 0E0N]{SP}, 代数空间的情形见~\cite[Tags 0E1D 和 0E1E]{SP}, 并见~\cite[定理 $4.17$]{DM69} (Deligne--Mumford 叠情形的陈述).
\end{proof}
还记录下面关于丰沛性的开放性的应用, 将命题~\bookref{B.2.10} 推广到代数空间.
\begin{proposition}[丰沛性的开放性]\label{src:5.6.17}
设 $X$ 是概形 $S$ 上的固有有限呈现代数空间, $L$ 是 $X$ 上的线丛. 若对某个几何点 $s:\Spec k\to S$, $L$ 在纤维 $X_s$ 上的限制 $L_s$ 丰沛, 则 $s$ 有开邻域 $U\subseteq S$, 使得 $X_U$ 是概形, 且 $L$ 在 $X_U$ 上的限制 $L_U$ 在 $U$ 上相对丰沛. 特别地, 对所有 $u\in U$, $L_u$ 在 $X_u$ 上丰沛.
\end{proposition}
\begin{proof}
应用 Gabber 引理 (\bookref{5.6.1}) 和丰沛性在有限覆盖下的不变性 (习题~\bookref{5.5.23}), 此陈述归结到概形的情形 (命题~\bookref{B.2.10}).
\end{proof}
